\documentclass{article}
\usepackage{amsmath}
\title{Adaptive Representation Exploration for Diverse Text Generation}
\author{Research Team}
\begin{document}
\maketitle

\begin{abstract}
Existing language-model decoding methods often generate semantically similar outputs. We present an adaptive representation exploration method that constructs a continuous behavioral space, samples candidate representations according to local density, and decodes multiple diverse outputs while preserving task validity. Experiments show improved diversity under comparable validity.
\end{abstract}

\section{Introduction}
Conventional sampling changes token probabilities but does not explicitly control high-level behavior. This causes mode concentration and repeated semantic patterns. The objective is to generate multiple valid outputs with controllable behavioral differences.

\section{Method}
Given an input prompt, the method first extracts a seed representation. It constructs a continuous behavioral space from concept anchors. Candidate representations are sampled according to a density-aware distribution. A radius calibration operation keeps candidates within a valid neighborhood. A diversity filter removes redundant candidates. The selected representations are inserted into a soft prompt and decoded to produce multiple outputs.

The calibrated representation is computed by
\begin{equation}
\hat{u}=u_0+\gamma(u-u_0).
\end{equation}

\section{Experiments}
Experiments are conducted on product design, role answering, and advertising generation tasks. The proposed method improves the pairwise diversity score while maintaining a comparable validity rate. Ablation results show that density-aware sampling and radius calibration both contribute to the final performance.

\begin{figure}
\caption{Overall pipeline of the continuous behavioral-space exploration method.}
\label{fig:pipeline}
\end{figure}

\section{Conclusion}
The method explores a continuous behavioral space and generates multiple valid outputs with different behaviors.
\end{document}

